\label{sec:implementation}

The \textsc{bisect-apportion} crate implements Jefferson/d'Hondt
using the priority formula $P_i^{\text{Jeff}}(s) = p_i/(s+1)$
(giving the $(s+1)$-th seat when already holding $s$ seats).
Starting from 0 seats per entity (not the constitutional-floor
initialisation used for HH, Webster, and Adams), the algorithm
awards all $H$ seats via priority queue and then applies the
constitutional floor $\max(1, \cdot)$ as a post-processing step.

\begin{verbatim}
bisect apportion --year 2020 --method jefferson
bisect apportion --year 2020 --method dhondt   # alias for jefferson
\end{verbatim}

\subsection{Integer Arithmetic}

The priority comparison for Jefferson is:
$P_i > P_j \iff p_i/(s_i+1) > p_j/(s_j+1) \iff p_i(s_j+1) > p_j(s_i+1)$.
This requires only 64-bit integer multiplication
(maximum value: $4 \times 10^7 \times 435 \approx 1.7 \times 10^{10}$,
well within \texttt{i64} range).

\subsection{D'Hondt Mode for PR Analysis}

The crate supports d'Hondt mode for proportional representation analysis:
\begin{verbatim}
bisect apportion --method dhondt --votes "PartyA:340,PartyB:280,..."
                 --seats 10
\end{verbatim}
This enables comparative analysis of how vote distributions translate
to seat allocations under d'Hondt vs.\ Sainte-Lagu\"e (Webster),
a common task in electoral systems research and policy analysis
for countries considering PR system reform.

\subsection{Test Suite}

\begin{itemize}
  \item \textbf{L0}: Priority formula $P(0) = p$, $P(1) = p/2$, $P(2) = p/3$.
  \item \textbf{L1}: 5-party, 10-seat d'Hondt example from
    \citet{lijphart1994electoral} produces the expected allocation
    [4, 3, 2, 1, 0] (with constitutional floor producing [4, 3, 2, 1, 1]
    in the U.S.\ apportionment context).
  \item \textbf{L2}: Jefferson for 2020 Census produces correct
    total of 435 seats with floor rounding applied.
\end{itemize}
